\chapter{Positive Matrices}\label{linear:ch:perron}

Suppose each adult pair of rabbits produces one juvenile pair per month,
juveniles become adults after one month, and adults survive. If $J_k,A_k$
count juvenile and adult pairs, the rule is
\[
\binom{J_{k+1}}{A_{k+1}}=
\begin{pmatrix}0&1\\1&1\end{pmatrix}\binom{J_k}{A_k}.
\]
Will the ratio of juveniles to adults settle down? More generally, when we
multiply a non-negative vector by a matrix repeatedly, does its direction
approach a limit? An eigenvalue larger in absolute value than all the others
would suggest an answer, but we must prove that such an eigenvalue exists.

Perron's theorem gives a positive eigenvalue whose modulus is strictly larger
than that of every other eigenvalue, and a corresponding eigenvector with
every coordinate positive. That eigenvalue is simple, so no Jordan block
can add a competing term at the same growth rate.

Zeros make the situation subtler. For a non-negative matrix, the directed graph of the non-zero entries matters. Cycles in that graph can create several eigenvalues on the outer spectral circle; their positions are governed by the greatest common divisor of cycle lengths.

This appendix develops the argument with small matrices, then connects it to
Markov chains, PageRank and population models. It uses the Jordan form proved
in Chapter~7 (Theorem~\ref{thm:jordan}) and the sequence limits proved in Volume~1,
Chapter~3. The finite-dimensional compactness step is supplied below.

\section{Positive matrices: Perron's theorem}

Throughout this appendix matrices are real and square of size $n\ge1$, unless
another domain is specified; eigenvectors may be complex. We write $A\ge0$ when all entries of $A$ (a matrix or a vector) are non-negative, $A>0$ when they are all positive, and
$|x|$ for the vector of moduli of the entries of $x$. The \emph{spectral radius} of a square matrix is $\rho(A)=\max|\lambda|$
over its complex eigenvalues. Two remarks are used constantly: if $A\ge0$, then $|Ax|\le A|x|$ (triangle inequality in each
coordinate); and if $A>0$, $x\ge0$ and $x\ne0$, then $Ax>0$.

For $A\ge0$ and $x\ge0$, $x\ne0$, the \emph{Collatz--Wielandt function} is
\[
r(x)=\min_{i:\ x_i>0}\frac{(Ax)_i}{x_i},
\]
the largest $t$ with $Ax\ge tx$; coordinates with $x_i=0$ impose no
restriction because $(Ax)_i\ge0$.

\begin{theorem}[Perron, 1907]\label{linear:thm:perron}
Let $A>0$ be an $n\times n$ matrix and $\rho=\rho(A)$. Then:
\begin{enumerate}[label=(\alph*)]
\item $\rho>0$ is an eigenvalue of $A$, with an eigenvector $z>0$ (the \emph{Perron vector});
\item $\rho$ is an algebraically simple eigenvalue;
\item every other eigenvalue satisfies $|\lambda|<\rho$;
\item every eigenvector $x\ge0$ of $A$ is a multiple of $z$.
\end{enumerate}
Moreover $\rho=\max_{x\ge0,x\ne0}r(x)$ (Collatz--Wielandt formula).
\end{theorem}
\begin{proof}
\emph{Existence and the variational formula.}
Let
$\Delta=\{x\in\mathbb R^n:x\ge0,\ \sum_i x_i=1\}$
and set $r^*=\sup_{x\in\Delta}r(x)$.
This supremum is finite: if $x_i$ is
a largest coordinate of $x$, then
$x_i>0$ and
\[
r(x)\le\frac{(Ax)_i}{x_i}
\le\sum_j a_{ij}
\le\max_i\sum_j a_{ij}.
\]
Whenever $Ax\ge t x$ with $t\ge0$,
positivity of $A$ gives $A(Ax)\ge tAx$.
In particular $r(Ax)\ge r(x)$,
and $r$ is unchanged by positive scaling.

The set
\[
K=\left\{\frac{Ax}{\sum_i(Ax)_i}:x\in\Delta\right\}
\]
stays away from the boundary of $\Delta$. Indeed, if $\alpha=\min_{ij}a_{ij}>0$
and $\beta=\max_{ij}a_{ij}$, then $\alpha\le(Ax)_i\le\beta$ for $x\in\Delta$,
so every coordinate of a point of $K$ is at least $\alpha/(n\beta)$.
Here is the compactness argument we need. From any sequence in $\Delta$,
apply the bounded-subsequence theorem of Volume~1, Chapter~3 to the first
coordinate, then the second, and so on. After finitely many extractions all
coordinates converge; their limits are non-negative and sum to $1$.
The map $x\mapsto Ax/\sum_i(Ax)_i$ is continuous, with a positive denominator,
so the same argument works for its image $K$. This is the finite
sequential compactness used here. On $K$, each ratio $(Ay)_i/y_i$ is continuous,
and so is their finite minimum: changing each ratio by at most $\varepsilon$
changes the minimum by at most $\varepsilon$. A sequence of points whose
$r$-values approach the supremum therefore has a subsequence attaining it
in the limit. Since $K\subseteq\Delta$
and $r(Ax)\ge r(x)$, we have
$\sup_K r=\sup_\Delta r=r^*$.
Choose $z\in K$ attaining this maximum.
Thus $z>0$ and $Az\ge r^*z$.
If $Az-r^*z$ were non-zero, it would be
a non-negative vector with at least
one positive coordinate. Applying
the strictly positive matrix $A$
would give $A(Az)>r^*Az$ in
\emph{every} coordinate, hence
$r(Az)>r^*$, an impossibility.
Therefore
\[
Az=r^*z,\qquad r^*>0.
\]

For any complex eigenpair $Ay=\lambda y$,
the triangle inequality gives
$A|y|\ge|\lambda||y|$.
By the definition of $r$ (also when
$|y|$ has zero coordinates),
$|\lambda|\le r(|y|)\le r^*$.
Thus the positive eigenvalue $r^*$
dominates all eigenvalue moduli:
$r^*=\rho$. This proves (a) and
the Collatz--Wielandt formula.

\emph{No other eigenvalue on the spectral circle.}
Suppose $Ay=\lambda y$, $y\ne0$ and
$|\lambda|=\rho$. Then
$A|y|\ge\rho|y|$. If this inequality
were not equality, applying $A$ and
the preceding strict-positivity
argument would give $r(A|y|)>\rho$,
contrary to maximality. Therefore
\[
A|y|=\rho|y|.
\]
Since $A|y|>0$, every coordinate
of $|y|$ is positive. In each row,
equality holds in
$\bigl|\sum_j a_{ij}y_j\bigr|
 \le\sum_j a_{ij}|y_j|$.
All $a_{ij}$ are positive, so equality forces all $y_j$ to have one
common argument. To see the equality case, rotate their non-zero sum to the
positive real axis: each summand has real part at most its modulus, and
equality of the sums forces equality for every summand. Thus
$y=e^{i\theta}|y|$ and
\[
\lambda y=Ay=e^{i\theta}A|y|
=\rho y.
\]
Hence $\lambda=\rho$, proving (c).

\emph{Geometric simplicity.}
Let $y$ be a real eigenvector for $\rho$,
and change its sign if necessary so that
some coordinate is positive. Define
\[
t=\min_{i:y_i>0}\frac{z_i}{y_i}>0,
\qquad u=z-ty.
\]
Then $u\ge0$ and at least one
coordinate of $u$ is zero.
But $Au=\rho u$. If $u\ne0$,
strict positivity of $A$ would
make $Au>0$, contradicting
the zero coordinate of $\rho u$.
Thus $u=0$ and $y$ is a multiple
of $z$. For a complex eigenvector
apply this separately to its real
and imaginary parts; the eigenspace
over $\mathbb C$ is also one-dimensional.

\emph{Algebraic simplicity.}
Apply the existence argument to
$A^{\mathsf T}>0$. It has the same
spectral radius, so there is a
positive vector $w$ with
$w^{\mathsf T}A=\rho w^{\mathsf T}$.
If $\rho$ had algebraic multiplicity
greater than one, geometric simplicity
would force a Jordan chain of length
at least two, with
$(A-\rho I)v=z$ for some $v$.
Multiplication by $w^{\mathsf T}$
would give
\[
0=w^{\mathsf T}(A-\rho I)v
=w^{\mathsf T}z>0,
\]
a contradiction. This proves (b).

\emph{Non-negative eigenvectors.}
Finally, if $Ax=\lambda x$ with
$x\ge0$, $x\ne0$, then
$w^{\mathsf T}x>0$ and
\[
\lambda w^{\mathsf T}x
=w^{\mathsf T}Ax
=\rho w^{\mathsf T}x.
\]
So $\lambda=\rho$, and geometric
simplicity makes $x$ a multiple
of $z$, proving (d).
\end{proof}

\paragraph{Why this works.}
For the positive matrix $A=\begin{pmatrix}2&1\\1&2\end{pmatrix}$, the eigenvalues are $3$ and $1$, with eigenvectors $(1,1)$ and $(1,-1)$. The first vector has positive coordinates; the second changes sign. Starting from $(1,0)$, repeated multiplication gives
\[
A^t\binom10=\frac{3^t}{2}\binom11+\frac12\binom1{-1}.
\]
The first term dominates, and the direction tends to $(1,1)$. Perron's theorem says that positivity of \emph{every entry} forces a similar dominant direction even when no simple formula gives the eigenvalues. The key is to show that a largest-modulus eigenvector can be chosen positive, and that neither an independent eigenvector nor a second eigenvalue of the same modulus can compete with it.

The vector $w>0$ with $w^{\mathsf T}A=\rho w^{\mathsf T}$ is the \emph{left Perron vector}; we normalise $w^{\mathsf T}z=1$.

\section{Contraction and the power method}

The convergence proof explains how the Perron vector wins: a positive matrix squeezes all positive directions towards one.
For positive vectors $x,y$, the ratios $x_i/y_i$ lie in an interval $[m(x/y),M(x/y)]$; we call its length the
\emph{oscillation} of $x/y$. It is zero exactly when $x$ is a multiple of $y$.

\begin{lemma}\label{linear:lem:oscillation}
Let $A>0$, and let $c=\min a_{ij}/a_{kj}$ over all $i,j,k$, so $0<c\le1$. For positive $x,y$,
\[
\begin{gathered}
[m(Ax/Ay),M(Ax/Ay)]\subseteq[m(x/y),M(x/y)],\\
M(Ax/Ay)-m(Ax/Ay)\le(1-c^2)\bigl(M(x/y)-m(x/y)\bigr).
\end{gathered}
\]
\end{lemma}
\begin{proof}
Put $m=m(x/y)$ and $M=M(x/y)$, so componentwise
$my\le x\le My$. Since every entry of $A$ is
positive, applying $A$ preserves these inequalities:
\[
mAy\le Ax\le MAy.
\]
Division by the positive coordinates of $Ay$ proves
the inclusion of intervals.

For the sharper estimate, let $q=A(My-x)\ge0$.
Because $A(x-my)\ge0$ as well,
\[
0\le q_i\le(M-m)(Ay)_i.
\]
Write $\beta_i=q_i/(Ay)_i$. By the definition of
$c$, for every $u\ge0$ and all row indices $i,k$,
\[
(Au)_i=\sum_j a_{ij}u_j
 \ge c\sum_j a_{kj}u_j=c(Au)_k.
\]
Use this first with $u=My-x$ to get $q_i\ge cq_k$,
and then with $u=y$ (interchanging $i$ and $k$)
to get $(Ay)_i\le c^{-1}(Ay)_k$. Hence
\[
\beta_i=\frac{q_i}{(Ay)_i}
 \ge c^2\frac{q_k}{(Ay)_k}=c^2\beta_k.
\]
Choose $i$ where $\beta$ is smallest and $k$
where it is largest. Since
$0\le\beta_k\le M-m$, this gives
\[
\max_i\beta_i-\min_i\beta_i
 \le(1-c^2)\max_i\beta_i
 \le(1-c^2)(M-m).
\]
Finally $(Ax)_i/(Ay)_i=M-\beta_i$, whose
oscillation is exactly that of $\beta$.
The bound includes the degenerate case $c=1$:
then all output ratios agree after one step.
\end{proof}

\begin{theorem}[convergence of the power method]\label{linear:thm:powerconv}
Let $A>0$ with Perron vectors $z,w$, $w^{\mathsf T}z=1$. For every $x\ge0$, $x\ne0$, $A^kx/\rho^k\to(w^{\mathsf T}x)\,z$, and
$A^k/\rho^k\to zw^{\mathsf T}$. The convergence is geometric.
\end{theorem}
\begin{proof}
First suppose $x>0$ and put
$u_k=A^kx/\rho^k$. Since $Az=\rho z$,
applying Lemma~\ref{linear:lem:oscillation} to
$A^kx$ and $\rho^kz$ shows that the smallest
and largest coordinates of $u_k/z$ form
nested intervals $[m_k,M_k]$ with
\[
0\le M_k-m_k\le(1-c^2)^k(M_0-m_0)\qquad(k\ge1).
\]
The positive left Perron vector satisfies
$w^{\mathsf T}A=\rho w^{\mathsf T}$ and
$w^{\mathsf T}z=1$. Consequently
\[
w^{\mathsf T}u_k=w^{\mathsf T}x
 =\sum_i(w_iz_i)\frac{(u_k)_i}{z_i}.
\]
The coefficients $w_iz_i$ are positive and sum
to $1$, so $\gamma=w^{\mathsf T}x$ lies in
\emph{every} interval $[m_k,M_k]$. Thus
\[
\left|\frac{(u_k)_i}{z_i}-\gamma\right|
 \le(1-c^2)^k(M_0-m_0)\qquad(k\ge1),
\]
a quantitative geometric convergence to
$(w^{\mathsf T}x)z$.
If $c=1$, the intervals already have length zero after the first step;
otherwise $0<1-c^2<1$.

If $x\ge0$ is non-zero but not strictly positive,
then $Ax>0$ because $A>0$. Apply the established
result to $Ax$ and observe that
\[
\frac{A^kx}{\rho^k}
=\frac1\rho\frac{A^{k-1}(Ax)}{\rho^{k-1}},
\qquad
w^{\mathsf T}(Ax)=\rho w^{\mathsf T}x.
\]
This proves the same limit, with a geometric
bound for $k\ge1$. In particular it applies to
each standard basis vector $e_j$. The limit
of column $j$ of $A^k/\rho^k$ is $w_jz$;
collecting the finitely many columns gives
$A^k/\rho^k\to zw^{\mathsf T}$, again
geometrically.
\end{proof}

So for any non-zero non-negative starting vector, $A^kx/\|A^kx\|$ converges to
$z/\|z\|$: this is the \emph{power method}, and
PageRank uses this iteration for the transpose of its stochastic matrix.
Here $\|\cdot\|$ is the Euclidean norm. The proof gives a geometric contraction bound $1-c^2$. The spectral picture is sharper but needs a qualification:
put $r=\max_{\lambda\ne\rho}|\lambda|/\rho$, taking $r=0$ if there are no other eigenvalues. Jordan form gives an error of order
$O(k^{m-1}r^k)$, where $m$ is the largest Jordan-block size among the eigenvalues of
modulus $r\rho$ (when $r>0$). In particular, a non-trivial Jordan block may add a
polynomial factor, so a uniform bound $O(r^k)$ does not always hold. For every
$r<q<1$ the error is nevertheless $O(q^k)$; if $r=0$, the non-Perron part vanishes
after finitely many powers. This is the precise sense in which the spectral gap
controls the asymptotic rate.
Indeed, a Jordan block $J=\lambda I+N$ with $N^e=0$ has
$J^k=\sum_{j=0}^{\min(k,e-1)}\binom{k}{j}\lambda^{k-j}N^j$.
For $\lambda\ne0$, its entries are bounded by a constant times
$k^{e-1}|\lambda|^k$; a zero-eigenvalue block vanishes for $k\ge e$.
If $0<a<1$, the ratio of successive terms of $k^d a^k$ tends to $a$,
so these terms eventually decrease geometrically and are bounded.
This absorbs all polynomial factors attached to smaller eigenvalue moduli,
and proves the stated $O(q^k)$ bound. When $r>0$, it also gives the sharper
$O(k^{m-1}r^k)$ bound from the blocks at modulus $r\rho$.

\begin{remark}
The oscillation measures direction rather than length: it vanishes precisely
when $x$ and $y$ lie on the same positive ray $\{tx:t>0\}$. Another measure is
$d(x,y)=\log(M(x/y)/m(x/y))$. It is unchanged by positive rescaling, symmetric,
and non-negative, with value zero exactly on one ray. For any positive $x,y,z$, the inequalities
$M(x/z)\le M(x/y)M(y/z)$ and $m(x/z)\ge m(x/y)m(y/z)$ give its triangle
inequality after taking logarithms. Thus it is a metric on positive rays,
called Hilbert's projective metric. The elementary oscillation estimate above
already provides every contraction bound used in this appendix; Birkhoff's
paper is listed below for further reading.
\end{remark}

\section{Non-negative matrices: Frobenius's theorem}

Zeros change the picture. The matrix $\bigl(\begin{smallmatrix}0&1\\1&0\end{smallmatrix}\bigr)$ has eigenvalues $\pm1$, both
of modulus $\rho$; the matrix $\bigl(\begin{smallmatrix}1&1\\0&1\end{smallmatrix}\bigr)$ has a Jordan block at $\rho$, and its
eigenvector $(1,0)$ is not positive. The right notion is connectivity. To $A\ge0$ of size $n$ associate the directed graph
$\Gamma(A)$ on vertices $1,\dots,n$ with an edge $i\to j$ when $a_{ij}>0$. Then $(A^k)_{ij}>0$ if and only if there is a walk of
length $k$ from $i$ to $j$, since $(A^k)_{ij}$ is a sum of products $a_{ii_1}a_{i_1i_2}\cdots a_{i_{k-1}j}$ over such walks.

\begin{definition}\label{linear:def:irreducible}
$A\ge0$ is \emph{irreducible} if $\Gamma(A)$ is strongly connected: there is a walk from every vertex to every vertex,
allowing a walk of length zero from a vertex to itself. For an irreducible
matrix with at least one edge, its \emph{period} $h\ge1$ is the greatest common
divisor of the \emph{positive} lengths of closed walks. It is \emph{primitive}
if $A^k>0$ for some integer $k\ge1$.
\end{definition}

The sole irreducible matrix without an edge is $[0]$: it has spectral radius
$0$, a simple eigenvalue with eigenvector $[1]$, no positive closed-walk length
and hence no period under our definition, and is not primitive. We handle it
separately. Every other irreducible graph has an outgoing edge at each vertex
and a positive closed walk. Delete a segment between repeated vertices in a
closed walk, repeating until only simple directed cycles remain. Each deletion
shortens a positive integer length, so this process terminates. Consequently
the gcd of closed-walk lengths equals the gcd of the finitely many simple
cycle lengths, each at most $n$. Bezout's identity therefore applies to this gcd.

Deleting repeated segments also gives a path of length at most $n-1$ between
any reachable vertices. Thus $A$ is irreducible if and only if $(I+A)^{n-1}>0$: the
$(i,j)$ entry of $(I+A)^{n-1}=\sum_k\binom{n-1}kA^k$ is positive exactly when some walk of length at most $n-1$ joins $i$ to $j$.

\begin{proposition}\label{linear:prop:levels}
Let $\Gamma$ be a finite strongly connected directed graph with at least one edge and period $h$. Fix a vertex $v_0$. The length of a walk from $v_0$ to $j$ is
well defined modulo $h$; calling it $\ell(j)$, every edge $i\to j$ satisfies $\ell(j)=\ell(i)+1$ in $\Z/h$. So the vertices
fall into $h$ classes, and every edge goes from one class to the next, cyclically.
\end{proposition}
\begin{proof}
Strong connectivity supplies a directed walk from
$v_0$ to any vertex $j$, and a directed walk
back from $j$ to $v_0$. Take two walks from
$v_0$ to $j$, with lengths $p$ and $q$, and
append the same return walk, of length $s$.
The resulting closed walks, both based at $v_0$,
have lengths $p+s$ and $q+s$.
By the definition of the period, $h$ divides
both lengths and therefore $p-q$. Thus the
residue class $\ell(j)$ of a walk length is
independent of the chosen walk.

If $i\to j$ is an edge, append it to any walk
from $v_0$ to $i$. This gives
$\ell(j)\equiv\ell(i)+1\pmod h$.
It remains to show that \emph{all} $h$ residue
classes really occur. A strongly connected
finite graph has at least one outgoing edge
at each vertex (including a loop if it has only
one vertex). Starting at $v_0$, follow an
outgoing edge repeatedly for $h-1$ steps.
The vertices visited after $0,1,\ldots,h-1$
steps have residues $0,1,\ldots,h-1$,
respectively, and therefore lie in distinct
classes. We have found all $h$ classes.
Every edge goes from one class to its cyclic
successor, as asserted.
\end{proof}

\begin{theorem}[Frobenius, 1912]\label{linear:thm:frobenius}
Let $A\ge0$ be non-zero and irreducible, of period $h$, and $\rho=\rho(A)$. Then $\rho>0$ is an algebraically simple
eigenvalue with an eigenvector $z>0$, and every eigenvector $x\ge0$ is a multiple of $z$. The eigenvalues of modulus $\rho$ are
exactly $\rho\,\omega^k$, $k=0,\dots,h-1$, where $\omega=e^{2\pi i/h}$; each is simple, and the whole spectrum is invariant under
multiplication by $\omega$.
\end{theorem}
\begin{proof}
For $A=[a]$ with $a>0$, the period is $1$ and all claims follow directly, so
assume $n\ge2$. Let $\Gamma$ be the
directed graph of positive entries.
By irreducibility it is strongly connected,
so between any two vertices there is
a path of length at most $n-1$.
The identity term in $I+A$ allows us to
pad such a path to length $n-1$;
therefore
\[
B=(I+A)^{n-1}>0.
\]
Moreover $A$ commutes with $B$.

\emph{The positive eigenvalue.}
Perron's theorem gives $B$ a positive
eigenvector $z$, spanning its unique
Perron eigenspace. Since $A$ commutes
with $B$, $Az$ belongs to this line:
$Az=rz$ for some $r\ge0$.
Apply the same argument to $A^{\mathsf T}$
to obtain $w>0$ with
$w^{\mathsf T}A=r'w^{\mathsf T}$.
From $w^{\mathsf T}Az$ we find
\[
r\,w^{\mathsf T}z=r'\,w^{\mathsf T}z,
\]
so $r'=r$. For any complex eigenpair
$Ay=\lambda y$,
\[
A|y|\ge|\lambda||y|.
\]
Multiply by $w^{\mathsf T}>0$.
Since $w^{\mathsf T}|y|>0$, we get
$|\lambda|\le r$. Thus $r=\rho(A)$.
If $r=0$, the equation $Az=0$ with
$z>0$ and $A\ge0$ forces $A=0$,
contrary to strong connectivity
when $n\ge2$. Hence $\rho>0$.

The eigenvalues of $B$ are
$(1+\lambda)^{n-1}$, counting
algebraic multiplicity, as $\lambda$
runs through the eigenvalues of $A$.
Furthermore
\[
|1+\lambda|\le1+|\lambda|\le1+\rho,
\]
with equality throughout only for
$\lambda=\rho$. The eigenvalue
$(1+\rho)^{n-1}$ of $B$ is
algebraically simple by Perron's theorem.
It follows that $\rho$ itself is
algebraically (and hence geometrically)
simple for $A$.
If $Ax=\lambda x$ and $x\ge0$ is non-zero,
the positive left eigenvector gives
$\lambda w^{\mathsf T}x
=\rho w^{\mathsf T}x$; thus
$\lambda=\rho$ and $x$ is a multiple
of $z$.

\emph{Which other eigenvalues can reach the
spectral circle?}
Let $Ay=\lambda y$ with
$|\lambda|=\rho$ and write
$\lambda=\rho e^{i\varphi}$.
The vector $A|y|-\rho|y|$ is
non-negative. Its scalar product with
$w>0$ vanishes, so it must be zero.
Thus $A|y|=\rho|y|$, and the
one-dimensional non-negative
eigenspace gives $|y|=cz>0$.
Write $y_j=|y_j|e^{i\theta_j}$.
In row $i$, equality in the triangle
inequality for
$\sum_j a_{ij}y_j
=\rho e^{i\varphi}y_i$ forces
\[
e^{i\theta_j}
=e^{i\varphi}e^{i\theta_i}
\qquad\text{whenever }a_{ij}>0.
\]
Along a closed directed walk of
length $\ell$, successive phases
therefore multiply by $e^{i\ell\varphi}$,
but the endpoint is the starting
vertex. Hence $e^{i\ell\varphi}=1$
for every closed walk length $\ell$.
Apply Bezout to the finitely many simple cycle lengths, whose gcd is $h$:
an integer combination of their lengths equals $h$, so $e^{ih\varphi}=1$.
Every peripheral eigenvalue is
therefore among $\rho\omega^k$,
where $\omega=e^{2\pi i/h}$.

\emph{Every permitted eigenvalue occurs.}
By Proposition~\ref{linear:prop:levels},
the vertices have labels
$\ell(i)\in\mathbb Z/h$ such that
$\ell(j)=\ell(i)+1$ whenever
$a_{ij}>0$. Put
$D=\operatorname{diag}
 (\omega^{\ell(1)},\ldots,\omega^{\ell(n)})$.
Then
\[
(D^{-1}AD)_{ij}
=\omega^{\ell(j)-\ell(i)}a_{ij}
=\omega a_{ij},
\]
including zero entries. Thus
$D^{-1}AD=\omega A$: multiplication
of the spectrum by $\omega$ preserves
eigenvalues with algebraic multiplicity.
Starting from the simple eigenvalue
$\rho$, we obtain all the
$\rho\omega^k$, each simple.
There are no others on the spectral
circle, completing the proof.
\end{proof}

\begin{figure}[ht]
\centering
\fig{linear-ch07-period}
\caption{A strongly connected digraph of period $3$. Every edge goes from one colour class to the next, so all closed walks
have lengths divisible by $3$; the cycle $1\to3\to6\to1$ has length $3$, making the gcd exactly $3$.
The adjacency matrix has peripheral eigenvalues $\rho$, $\rho\omega$, $\rho\omega^2$, where $\omega=e^{2\pi i/3}$.}
\label{linear:fig:period}
\end{figure}

\section{Primitive matrices, Wielandt's bound}

\begin{theorem}\label{linear:thm:primitive}
For non-zero irreducible $A\ge0$, the following are equivalent: (i) $A$ is primitive; (ii) the period is $h=1$; (iii) $\rho$ is the
only eigenvalue of modulus $\rho$. For a primitive $A$, $A^k/\rho^k\to zw^{\mathsf T}$, where $z,w$ are the right and left Perron
vectors with $w^{\mathsf T}z=1$.
\end{theorem}
\begin{proof}
The equivalence of (ii) and (iii)
follows immediately from
Theorem~\ref{linear:thm:frobenius},
which lists the peripheral eigenvalues
as $\rho\omega^j$, $0\le j<h$.

\emph{Primitive implies one peripheral
eigenvalue.}
Suppose $A^q>0$ for some $q\ge1$.
Perron's theorem applied to $A^q$
says that $\rho^q$ is its sole
eigenvalue of modulus $\rho^q$
and is algebraically simple.
If $A$ had a different eigenvalue
$\lambda$ with $|\lambda|=\rho$,
then $\lambda^q$ would have modulus
$\rho^q$, hence equal $\rho^q$.
Independent eigenvectors of $A$
for $\rho$ and $\lambda$ would
be independent eigenvectors of
$A^q$ for $\rho^q$, contradicting
simplicity. Thus (i) implies (iii).

\emph{Period one gives positive
powers.}
Assume $h=1$ and fix a vertex $i$
of the directed graph $\Gamma(A)$.
Let $S$ be the positive lengths
of closed walks starting at $i$.
Concatenation shows that $S$ is
closed under addition.
Its greatest common divisor is
one. Indeed, given a closed walk
of length $\ell$ at any vertex $j$,
choose a path from $i$ to $j$
of length $a$ and one back of
length $b$. Closed walks at $i$
of lengths $a+b$ and $a+b+\ell$
show that $\gcd S$ divides $\ell$.
Because the gcd of all closed-walk
lengths is $h=1$, we have
$\gcd S=1$.

Choose finitely many lengths $s_1,\ldots,s_t\in S$ with gcd one.
For justification, start with any length and adjoin another whenever the
current gcd fails to divide some member of $S$. Each such step strictly
decreases a positive integer, so the process terminates at $\gcd S=1$.
Bezout's identity gives integers
$c_j$ with $\sum_jc_js_j=1$.
Group the positive and negative
terms to write $P-N=1$, where
$P,N$ are non-negative integer
combinations of lengths in $S$.
Hence $P,N\in S\cup\{0\}$.
If $N=0$, then $1=P\in S$ and
every positive length belongs to $S$.
Otherwise $N\ge1$ and
$P=N+1$. For every $m\ge N^2$,
write $m=qN+r$ with $0\le r<N$,
so $q\ge N>r$. Then
\[
m=rP+(q-r)N,
\]
a sum of closed-walk lengths.
Thus, in either case, all
sufficiently large integers lie
in $S$.

Fix a path from $i$ to any vertex
$j$, of length $b_{ij}$. Prefixing
that path by a sufficiently long
closed walk at $i$ produces
a walk from $i$ to $j$
of \emph{every sufficiently large
length}. There are finitely many
ordered pairs $(i,j)$, so one
integer $q$ works for them all.
Thus $(A^q)_{ij}>0$ for all
$i,j$, proving (ii) implies (i).

\emph{Convergence.}
Choose $q$ with $B=A^q>0$.
The Perron vectors $z,w$ of
$A$ are also those of $B$,
with eigenvalue $\rho^q$,
and normalise $w^{\mathsf T}z=1$.
Write $k=mq+s$, $0\le s<q$.
Theorem~\ref{linear:thm:powerconv}
applied to $B$ gives
\[
\frac{A^k}{\rho^k}
=\frac{B^m}{(\rho^q)^m}
  \frac{A^s}{\rho^s}
\longrightarrow
zw^{\mathsf T}\frac{A^s}{\rho^s}
=zw^{\mathsf T}.
\]
There are only finitely many
possible remainders $s$, so
this proves convergence along
the whole sequence.
\end{proof}

How large must $k$ be? Wielandt stated the sharp bound in 1950 without a
proof. Holladay and Varga published a proof in 1958. Schneider's 2002 paper
reproduces a proof found in Wielandt's mathematical diaries.

\begin{theorem}[Wielandt]\label{linear:thm:wielandt}
If $A\ge0$ of size $n\ge2$ is primitive, then $A^{(n-1)^2+1}>0$.
\end{theorem}
\begin{proof}
Consider the directed graph $\Gamma$ of positive entries of
$A$. A positive power of $A$ provides a walk between every ordered pair,
so $\Gamma$ is strongly connected. Theorem~\ref{linear:thm:primitive}
then gives period one.
Let $C$ be a shortest directed cycle, of length $s$.
A simple directed cycle has at most $n$ vertices, so $s\le n$.
If $s=n$, every directed cycle must have length $n$:
a shorter one would contradict minimality, and a longer
simple cycle is impossible. Every closed walk decomposes
into directed cycles, so all closed-walk lengths would be
multiples of $n$, contradicting period one for $n\ge2$.
Therefore
\[
1\le s\le n-1.
\]

Set $B=A^s$. Every sufficiently large power of $A$
is positive, and consequently $B^k=A^{sk}>0$
for every sufficiently large $k$. Thus $B$ is primitive,
and in particular its directed graph is strongly connected.
At every vertex $c$ of $C$, following $C$ once gives a
closed walk of length $s$; hence $B_{cc}>0$, a loop in
the graph of $B$.

Fix $c\in C$ and any target vertex $j$.
Strong connectivity of the graph of $B$ supplies a
directed path from $c$ to $j$ of length at most $n-1$
(a shortest path visits no vertex twice).
Prepend enough copies of the loop at $c$ to make
its length exactly $n-1$. Therefore
\[
(B^{n-1})_{cj}=(A^{s(n-1)})_{cj}>0
\qquad(c\in C,\ j\in\Gamma).
\]

Now fix an arbitrary starting vertex $i$. Choose a
shortest directed path from $i$ to the vertex set $C$.
Before its first arrival at $C$ it visits distinct
vertices outside $C$, so its length $t$ is at most
$n-s$. Continue along the cycle $C$ for another
$n-s-t$ steps. We reach some $c\in C$ after
exactly $n-s$ steps, which gives
$(A^{n-s})_{ic}>0$. Following this by the
length-$s(n-1)$ walk from $c$ to any $j$ proves
\[
(A^{n-s+s(n-1)})_{ij}
\ge (A^{n-s})_{ic}(A^{s(n-1)})_{cj}>0.
\]
Hence $A^N>0$ for $N=n+s(n-2)$.

Since $s\le n-1$ and $n\ge2$,
\[
N=n+s(n-2)
\le n+(n-1)(n-2)=(n-1)^2+1.
\]
Finally, a primitive irreducible matrix has no zero
row; whenever $A^k>0$, the product $A^{k+1}=AA^k$
is also positive. Repeating this observation
extends positivity from exponent $N$ to
$(n-1)^2+1$, as claimed.
\end{proof}

The bound is attained by the \emph{Wielandt matrix}, the graph
$1\to2\to\dots\to n\to1$ with one extra edge $n\to2$.
Exercise~\ref{linear:ex:wielandtsharp} proves sharpness for every $n\ge2$.

\section{Markov chains and PageRank}

A \emph{stochastic} matrix $P\ge0$ has rows summing to $1$: $P\mathbf 1=\mathbf 1$. It describes a Markov chain, $p_{ij}$ being the
probability to move from state $i$ to state $j$, using the same rule at each
step and depending only on the current state. For a finite path
$i_0,\ldots,i_k$, assign probability
$\mu_{i_0}p_{i_0i_1}\cdots p_{i_{k-1}i_k}$, where $\mu$ is the initial
probability row. Summing successively over the last state shows that these
probabilities sum to $1$. Summing over paths ending at $j$ gives the $j$-th
entry of $\mu P^k$. In particular one step sends $\mu$ to $\mu P$. The vector $\mathbf 1>0$ is a non-negative eigenvector for the eigenvalue $1$, so $\rho(P)=1$ when $P$ is irreducible
(Theorem~\ref{linear:thm:frobenius}: non-negative eigenvectors belong to $\rho$), and indeed for every stochastic $P$, since $\|Px\|_\infty
\le\|x\|_\infty$, where $\|x\|_\infty=\max_i|x_i|$.

\begin{corollary}[ergodic theorem for Markov chains]\label{linear:cor:markov}
Let $P$ be a stochastic matrix that is irreducible with period $1$. There is a unique probability vector $\pi$ with $\pi P=\pi$;
it is positive, and for every initial distribution $\mu$, $\mu P^k\to\pi$. More precisely $P^k\to\mathbf 1\pi$, the matrix all of
whose rows are $\pi$.
\end{corollary}
\begin{proof}
We regard distributions as row vectors.
Since $P$ is row-stochastic,
$P\mathbf 1=\mathbf 1$. In particular $1$
is an eigenvalue, and the maximum-row-sum
bound gives $\rho(P)\le1$, so $\rho(P)=1$.

Irreducibility gives a strictly positive left
Perron eigenvector by
Theorem~\ref{linear:thm:frobenius}. Normalising
its coordinates to sum to $1$ produces a
probability row vector $\pi$ satisfying
$\pi P=\pi$. The Perron eigenspace is
one-dimensional, so this stationary probability
vector is unique. Because the period is $1$,
Theorem~\ref{linear:thm:primitive} applies and
all other eigenvalues have modulus less than
$1$. Its convergence statement, with right
Perron vector $\mathbf 1$ and left vector
$\pi$, reads
\[
P^k\longrightarrow\mathbf1\pi.
\]
For any initial probability row vector $\mu$,
we have $\mu\mathbf1=1$, and hence
\[
\mu P^k\longrightarrow
\mu(\mathbf1\pi)=(\mu\mathbf1)\pi=\pi.
\]
The assertion requires aperiodicity:
irreducibility alone would permit oscillating
distributions, as in the two-state alternating
chain.
\end{proof}

Volume~4, Chapter~9 returns to convergence by coupling: it compares two chains
directly. Here the relevant number is the largest modulus $r$ of an eigenvalue
other than $1$, or $r=0$ if none exists. The \emph{absolute spectral gap} is
$1-r$. The Jordan-block calculation above gives error $O(q^k)$ for every
$r<q<1$; a bound $O(r^k)$ needs further hypotheses, such as diagonalizability.

\begin{example}[PageRank]\label{linear:exa:pagerank}
Model a surfer on the web: from page $i$ with $d_i$ outgoing links, follow one at random ($p_{ij}=1/d_i$ for each link
$i\to j$; a page with no links is given links to all pages). We seek an equilibrium probability for each page, a stationary distribution
of $P$. Such a web graph need not be strongly connected. Add teleportation,
as in the PageRank model: with probability $1-\alpha$, the surfer jumps to
a uniformly random page. The \emph{Google matrix}
\[
G=\alpha P+\frac{1-\alpha}n\mathbf 1\mathbf 1^{\mathsf T},\qquad\alpha=0.85,
\]
is stochastic and positive, so Perron's theorem gives a unique positive PageRank vector $\pi=\pi G$, and the power method
has an error bound independent of the size of the web. For a real row $y$,
write $\|y\|_1=\sum_i|y_i|$. Stochasticity gives
\[
\|yP\|_1\le\sum_j\sum_i|y_i|p_{ij}=\|y\|_1.
\]
If $y\mathbf1=0$, then $yG=\alpha yP$, because the teleportation term vanishes.
Taking $y=\mu-\pi$ and iterating gives
$\|\mu G^k-\pi\|_1\le2\alpha^k$ for any initial probability row $\mu$.
The coordinates of the difference sum to zero, so its positive coordinates
sum to half its $1$-norm, as do the absolute values of its negative coordinates.
Thus the error in each coordinate is at most $\alpha^k$.
At $\alpha=0.85$, $50$ steps guarantee only an error below $0.000296$ per
coordinate; $86$ steps guarantee an error below $10^{-6}$.
Figure~\ref{linear:fig:pagerank} shows a small example. Page $C$, with six incoming links, ranks first. Page $A$
has only two, but one of them comes from $C$, which has no other outgoing link
and passes its entire link-following share to $A$: so $A$ ranks
almost as high as $C$. Page $B$ also has two incoming links, one from the important page $A$, and ranks third. Being linked by
important pages matters more than being linked often.
\end{example}

\begin{figure}[ht]
\centering
\fig{linear-ch07-pagerank}
\caption{PageRank of an eight-page web with $\alpha=0.85$; the area of each node is approximately proportional to its rank. The values
are $C:0.351$, $A:0.329$, $B:0.169$, and between $0.019$ and $0.040$ for the others.}
\label{linear:fig:pagerank}
\end{figure}

\section{Populations}

\begin{example}[Leslie models]\label{linear:exa:leslie}
Divide a female population into age classes $1,\dots,n$. Each step, a female of class $i$ gives birth on average to $f_i\ge0$
daughters, and for $i<n$ survives into class $i+1$ with probability $s_i\in(0,1]$. The resulting mean-population model is
$x_{k+1}=Lx_k$, where $x_k$ records mean class sizes and the \emph{Leslie matrix} is
\[
L=\begin{pmatrix}f_1&f_2&\cdots&f_{n-1}&f_n\\s_1&0&\cdots&0&0\\0&s_2&&&0\\&&\ddots&&\vdots\\0&0&\cdots&s_{n-1}&0\end{pmatrix}.
\]
If $f_n>0$, the graph $\Gamma(L)$ contains the cycle $1\to n\to n-1\to\dots\to2\to1$ (the edge $1\to j$ comes from
$f_j>0$ and the edge $j+1\to j$ from $s_j>0$), so $L$ is irreducible; if moreover two consecutive $f_i,f_{i+1}$ are positive,
there are cycles $1\to i\to\dots\to1$ of lengths $i$ and $i+1$, the period is $1$ and $L$ is primitive.
Whenever $L$ is primitive, for any $x_0\ge0$, $x_0\ne0$, and
$w^{\mathsf T}z=1$, we have
$L^kx_0/\rho^k\to(w^{\mathsf T}x_0)z$ with $w^{\mathsf T}x_0>0$.
The total population's successive ratio tends to $\rho$, and the class
proportions tend to the \emph{stable age distribution} $z/\sum_i z_i$.
The survival rows of $Lz=\rho z$ give
$z_i=\ell_i\rho^{-(i-1)}z_1$, with $\ell_1=1$ and
$\ell_i=s_1\cdots s_{i-1}$ the probability to reach class $i$.
Substitute into the first row and divide by $\rho z_1>0$ to get the
\emph{Euler--Lotka equation} $\sum_i f_i\ell_i\rho^{-i}=1$.
Exercise~\ref{linear:ex:leslie} finds the full characteristic polynomial,
including reducible cases and zero eigenvalues.

Fibonacci's rabbits, counted in pairs, give a two-stage population model,
juveniles and adults, with $f_1=0$,
$f_2=1$, $s_1=1$, and one extra entry $1$ in position $(2,2)$ because adults remain adults:
$L=\bigl(\begin{smallmatrix}0&1\\1&1\end{smallmatrix}\bigr)$ in the coordinates $(J,A)$ of the opening question, with
$\rho=\varphi=\frac{1+\sqrt5}2$ and Perron vector $z=(1,\varphi)$ (indeed $Lz=(\varphi,1+\varphi)=\varphi z$). This $L$ is
primitive ($L^2>0$), so the fraction of juveniles tends to $\frac1{1+\varphi}=\frac1{\varphi^2}\approx0.382$, whatever the
nonzero non-negative start: the golden ratio appears both as the growth
factor and in the stable proportions.
\end{example}

\begin{beauty}
The rabbit matrix has a growth factor $\varphi$ and an adult-to-juvenile ratio
$\varphi$. Those are two readings of one eigenvector equation. PageRank uses
the left eigenvector instead: its coordinates record proportions rather than
counts. In both examples the graph tells us whether repeated multiplication
settles to one direction or keeps cycling between several classes.
\end{beauty}

\section{Exercises}

\begin{exercise}[$\star$]\label{linear:ex:perron2}
Let $A=\bigl(\begin{smallmatrix}a&b\\c&d\end{smallmatrix}\bigr)>0$. Compute $\rho$ and the Perron vector, and check directly
that the other eigenvalue has smaller modulus.
\end{exercise}

\begin{exercise}[$\star$]\label{linear:ex:rowsum}
For $A\ge0$ irreducible, show that $\min_i\sum_ja_{ij}\le\rho\le\max_i\sum_ja_{ij}$, with equality on either side only if all row
sums are equal.
\end{exercise}

\begin{exercise}[$\star\star$]\label{linear:ex:minmaxcw}
For $A>0$, show the dual Collatz--Wielandt formula $\rho=\min_{x>0}\max_i\frac{(Ax)_i}{x_i}$, and deduce that any $x>0$ gives
$\min_i\frac{(Ax)_i}{x_i}\le\rho\le\max_i\frac{(Ax)_i}{x_i}$.
\end{exercise}

\begin{exercise}[$\star\star$]\label{linear:ex:monotone}
Let $0\le A\le B$ entrywise, with $B$ irreducible. Show that $\rho(A)\le\rho(B)$, with equality only if $A=B$.
\end{exercise}

\begin{exercise}[$\star\star$]\label{linear:ex:nonneg}
Let $A\ge0$ be arbitrary. Show that $\rho(A)$ is an eigenvalue of $A$ with an eigenvector $x\ge0$, $x\ne0$. (Apply
Exercise~\ref{linear:ex:monotone} to $A+\varepsilon J$, where $J$ is the all-ones matrix, and let $\varepsilon\to0$.)
\end{exercise}

\begin{exercise}[$\star\star$]\label{linear:ex:cycleperiod}
Find the period and eigenvalues of the directed cycle on $n\ge1$ vertices.
For a connected bipartite graph with at least one edge, put its adjacency
matrix in the form $\bigl(\begin{smallmatrix}0&B\\B^{\mathsf T}&0\end{smallmatrix}\bigr)$.
Find its period and describe its full spectrum using the positive eigenvalues
of $BB^{\mathsf T}$ and the rank of $B$. In particular, prove that its spectrum,
with multiplicities, is symmetric about $0$.
\end{exercise}

\begin{exercise}[$\star$]\label{linear:ex:diagpos}
Show that an irreducible $A\ge0$ with a positive diagonal entry is primitive, and that if all diagonal entries are positive then
$A^{n-1}>0$.
\end{exercise}

\begin{exercise}[$\star$]\label{linear:ex:twostate}
For the two-state chain $P=\bigl(\begin{smallmatrix}1-a&a\\b&1-b\end{smallmatrix}\bigr)$ with $0<a,b<1$, find $\pi$ and an
explicit formula for $P^k$. What is the rate of convergence?
\end{exercise}

\begin{exercise}[$\star\star$]\label{linear:ex:pagerankrate}
Let $0\le\alpha<1$, let $P$ be stochastic, and let $v$ be a probability
column vector, possibly with zero coordinates. Show that every eigenvalue
$\lambda\ne1$ of $G=\alpha P+(1-\alpha)\mathbf1v^{\mathsf T}$ satisfies
$|\lambda|\le\alpha$. Prove that $G$ has a unique stationary probability
row $\pi$, and that $\|\mu G^k-\pi\|_1\le2\alpha^k$ for every probability
row $\mu$ and $k\ge1$, even when $G$ is not irreducible.
\end{exercise}

\begin{exercise}[$\star\star$]\label{linear:ex:leslie}
For the Leslie matrix with $n\ge1$, $f_i\ge0$ and $s_i>0$, set $\ell_1=1$.
Show that its characteristic polynomial is $X^n-\sum_if_i\ell_iX^{n-i}$, that it has exactly one
positive root when some $f_i>0$, and that this root is $\rho(L)$, including when $f_n=0$.
Compute a positive right Perron vector, and show that $\rho>1$ if and only
if the \emph{net reproductive rate} $\sum_if_i\ell_i$ exceeds $1$.
\end{exercise}

\begin{exercise}[$\star\star$]\label{linear:ex:doublystoch}
Let $P$ be an $n\times n$ primitive matrix with both row and column sums $1$
(\emph{doubly stochastic}). Show that $P^k\to J/n$.
A shuffle of $d$ distinct cards chooses a permutation of positions from a fixed
probability distribution, independently of the current deck order and of
earlier choices. Show that its transition matrix on the $d!$ orders is doubly
stochastic. Deduce convergence to the uniform distribution when that matrix
is primitive.
\end{exercise}

\begin{exercise}[$\star\star\star$]\label{linear:ex:wielandtsharp}
Let $n\ge2$ and let $W$ be the Wielandt matrix (edges $i\to i+1$ for $i<n$, $n\to1$ and $n\to2$). Show that every positive-length closed walk at vertex $1$
has length $an+b(n-1)$ with $a\ge1$ and $b\ge0$, and deduce that $(W^{(n-1)^2})_{11}=0$, so that Theorem~\ref{linear:thm:wielandt} is
sharp.
\end{exercise}

\section{Hints}

\begin{hint}{linear:ex:perron2}
Compute the discriminant of $(X-a)(X-d)-bc$. Which root exceeds both $a$ and $d$?
\end{hint}
\begin{hint}{linear:ex:rowsum}
Look at the row of $Az=\rho z$ where $z_i$ is largest, and the one where it is smallest.
\end{hint}
\begin{hint}{linear:ex:minmaxcw}
For $x>0$, apply the left Perron vector to $Ax\le Mx$.
\end{hint}
\begin{hint}{linear:ex:monotone}
If $Ay=\lambda y$ then $|\lambda||y|\le B|y|$; multiply by the left Perron vector of $B$. For equality, look at what is forced.
\end{hint}
\begin{hint}{linear:ex:nonneg}
$\rho(A+\varepsilon J)$ decreases as $\varepsilon\downarrow0$; take a limit of normalised Perron vectors.
\end{hint}
\begin{hint}{linear:ex:cycleperiod}
For a bipartite graph, a two-step return determines the period. Use the blocks
$B,B^{\mathsf T}$ to turn an eigenvector of $BB^{\mathsf T}$ into two
eigenvectors of the adjacency matrix; treat zero eigenvalues separately.
\end{hint}
\begin{hint}{linear:ex:diagpos}
A loop at a vertex makes $1$ a closed walk length. With all loops, walks can be padded to any longer length.
\end{hint}
\begin{hint}{linear:ex:twostate}
The eigenvalues are $1$ and $1-a-b$.
\end{hint}
\begin{hint}{linear:ex:pagerankrate}
A left eigenvector $x$ for $\lambda\ne1$ satisfies $x^{\mathsf T}\mathbf 1=0$.
On zero-sum rows, $yG=\alpha yP$ and $\|yP\|_1\le\|y\|_1$.
Bound the differences between consecutive iterates by a geometric series.
\end{hint}
\begin{hint}{linear:ex:leslie}
Expand along the first row, then solve the survival rows successively.
The last positive fertility defines an irreducible leading block, giving
existence of a positive root. Monotonicity of $\sum_i f_i\ell_i\lambda^{-i}$
gives uniqueness.
\end{hint}
\begin{hint}{linear:ex:doublystoch}
$\mathbf 1$ is both a left and a right Perron vector.
\end{hint}
\begin{hint}{linear:ex:wielandtsharp}
Vertex $n$ is the only vertex with two outgoing edges. Reduce $(n-1)^2=an+b(n-1)$ modulo $n-1$.
\end{hint}

\section{Solutions}

\begin{solution}{linear:ex:perron2}
\idea There are only two eigenvalues, and a positive matrix has a positive leading eigenvector. Solve the quadratic, then use the positivity of the trace to compare the moduli.

\writeup Let $A=\begin{pmatrix}a&b\\c&d\end{pmatrix}$ with $a,b,c,d>0$. Its characteristic polynomial is
\[
\chi_A(X)=X^2-(a+d)X+(ad-bc).
\]
The discriminant is
\[
\Delta=(a+d)^2-4(ad-bc)=(a-d)^2+4bc>0.
\]
Thus the eigenvalues are real and distinct:
\[
\lambda_\pm=\frac{a+d\pm\sqrt{(a-d)^2+4bc}}2.
\]
Because $bc>0$, the square root is strictly larger than $|a-d|$, hence
\[
\lambda_+>\frac{a+d+|a-d|}2=\max(a,d)>0.
\]
The vector
\[
z=\begin{pmatrix}b\\\lambda_+-a\end{pmatrix}
\]
has two positive coordinates. Its first eigenvalue equation is $(a-\lambda_+)b+b(\lambda_+-a)=0$, and the second follows from the characteristic equation:
\[
cb+(d-\lambda_+)(\lambda_+-a)
=bc-(\lambda_+-d)(\lambda_+-a)=0.
\]
Therefore $Az=\lambda_+z$.

Finally $\lambda_++\lambda_-=a+d>0$ and $\lambda_+>\lambda_-$. If $\lambda_-\ge0$ this immediately gives $|\lambda_-|<\lambda_+$. If $\lambda_-<0$, the trace identity gives $\lambda_+>-\,\lambda_-$, again $|\lambda_-|<\lambda_+$. Hence $\rho(A)=\lambda_+$, and $z$ is a Perron eigenvector, as claimed.
\end{solution}

\begin{solution}{linear:ex:rowsum}
\idea Compare $Az=\rho z$ with the row sums at a coordinate where a positive Perron eigenvector $z$ is largest, then smallest. For equality, keep track of when the inequalities can be equalities and use irreducibility to spread this equality to every coordinate.

\writeup For $A=[0]$, both row-sum bounds are equalities and the conclusion is immediate.
Otherwise let $A=(a_{ij})$ be irreducible with $a_{ij}\ge0$, and let $z>0$ be a right Perron eigenvector, so
\[
Az=\rho z,\qquad z_i>0\quad\text{for every }i.
\]
Write $r_i=\sum_j a_{ij}$ for the $i$-th row sum, and let $M=\max_i r_i$, $m=\min_i r_i$.

Choose an index $i$ with $z_i=\max_j z_j$. Since $a_{ij}\ge0$,
\[
\rho z_i=\sum_j a_{ij}z_j
\le\sum_j a_{ij}z_i=r_i z_i\le Mz_i.
\]
Division by the positive number $z_i$ gives $\rho\le M$.

At an index $k$ with $z_k=\min_j z_j$, reverse the first inequality:
\[
\rho z_k=\sum_j a_{kj}z_j
\ge\sum_j a_{kj}z_k=r_kz_k\ge mz_k.
\]
Thus $m\le\rho\le M$.

Suppose now that $\rho=M$. At an index $i$ of maximum $z_i$, equality must hold throughout
\[
Mz_i=\sum_j a_{ij}z_j\le r_i z_i\le Mz_i.
\]
Therefore $r_i=M$ and $a_{ij}>0$ forces $z_j=z_i$. Each such neighbour $j$ is also a maximum coordinate. Repeating the same argument from that neighbour forces its row sum to be $M$ and all its outgoing positive neighbours to have the same coordinate. Irreducibility of $A$ means that the directed graph with arrows $i\to j$ when $a_{ij}>0$ is strongly connected, so this propagation reaches every index. Hence $z$ is constant and
\[
A\mathbf1=M\mathbf1;
\]
all row sums equal $M$.

If instead $\rho=m$, start at a minimum coordinate and use equality in the lower-bound chain. The same propagation shows all row sums equal $m$.

Conversely, if all row sums equal a number $c$, then $A\mathbf1=c\mathbf1$. The exercise assumes $A$ irreducible and nonnegative, so $\mathbf1>0$ is a nonnegative eigenvector. Frobenius's theorem (Theorem~\ref{linear:thm:frobenius}) puts every such eigenvector on the spectral radius, and $c=\rho$. Equality holds on both sides. The equality characterisation is therefore exact.
\end{solution}

\begin{solution}{linear:ex:minmaxcw}
\idea A left Perron eigenvector turns a coordinatewise inequality into a numerical inequality. This is the simplest way to prove both sides of the Collatz--Wielandt estimate.

\writeup Let $A>0$. Perron's theorem provides right and left eigenvectors $z>0$ and $w>0$ with
\[
Az=\rho z,\qquad w^{\mathsf T}A=\rho w^{\mathsf T}.
\]
Take any $x>0$ and define
\[
m(x)=\min_i\frac{(Ax)_i}{x_i},
\qquad
M(x)=\max_i\frac{(Ax)_i}{x_i}.
\]
All ratios are defined because each $x_i$ is positive. By definition,
\[
m(x)x\le Ax\le M(x)x
\]
coordinate by coordinate. Multiply by the positive row vector $w^{\mathsf T}$:
\[
m(x)\,w^{\mathsf T}x
\le w^{\mathsf T}Ax
=\rho\,w^{\mathsf T}x
\le M(x)\,w^{\mathsf T}x.
\]
The scalar $w^{\mathsf T}x$ is strictly positive, so division gives
\[
\boxed{m(x)\le\rho\le M(x)}
\]
for every $x>0$.

In particular $\rho\le\inf_{x>0}M(x)$. But for the positive Perron vector $x=z$, every ratio equals $\rho$:
\[
\frac{(Az)_i}{z_i}=\rho.
\]
Therefore $M(z)=\rho$, so the infimum is attained at $z$ and equals $\rho$. Thus
\[
\boxed{\rho=\min_{x>0}\max_i\frac{(Ax)_i}{x_i}}.
\]
The same argument gives the dual max--min formula. A test vector need not be an eigenvector to produce useful bounds; only at the Perron vector do the bounds coincide.
\end{solution}

\begin{solution}{linear:ex:monotone}
\idea Increasing matrix entries cannot decrease the spectral radius of a non-negative matrix. To prove strictness when the larger matrix is irreducible, use its strictly positive Perron vectors: a non-zero non-negative difference then contributes a strictly positive scalar.

\writeup Assume $0\le A\le B$ entrywise and $B$ irreducible.
If $B=[0]$, then $A=[0]$ too and the assertion is immediate.
Otherwise let $w>0$ be a left Perron eigenvector of $B$:
\[
w^{\mathsf T}B=\rho(B)w^{\mathsf T}.
\]
Choose an eigenvalue $\lambda$ of $A$ with $|\lambda|=\rho(A)$, and a non-zero complex eigenvector $y$. Taking absolute values componentwise gives
\[
|\lambda|\,|y|=|Ay|\le A|y|\le B|y|.
\]
Multiply by $w^{\mathsf T}$ and use its eigenvector identity:
\[
\rho(A)w^{\mathsf T}|y|
\le w^{\mathsf T}B|y|
=\rho(B)w^{\mathsf T}|y|.
\]
Since $w>0$ and $|y|\ne0$, the common scalar factor is positive, proving
\[
\boxed{\rho(A)\le\rho(B)}.
\]

Suppose equality holds. Every term of the chain is now equal after multiplication by the positive vector $w$. In particular
\[
w^{\mathsf T}\bigl(B|y|-\rho(B)|y|\bigr)=0,
\]
where the vector in parentheses is componentwise non-negative. Since $w>0$, this vector must vanish, giving
\[
B|y|=\rho(B)|y|.
\]
By irreducibility, Frobenius's theorem (Theorem~\ref{linear:thm:frobenius}) makes a non-zero non-negative eigenvector strictly positive, so $|y|>0$.

Equality also gives
\[
w^{\mathsf T}(B-A)|y|=0.
\]
Every entry of $B-A$, $w$, and $|y|$ is non-negative, with the latter two strictly positive. If any entry of $B-A$ were positive, the scalar above would be positive. Hence $B-A=0$, or $A=B$. Conversely $A=B$ certainly gives equality. Thus
\[
\boxed{A\ne B\implies\rho(A)<\rho(B)}
\]
under the irreducibility hypothesis on $B$.
\end{solution}

\begin{solution}{linear:ex:nonneg}
\idea Approximate a non-negative matrix by strictly positive ones. Each approximation has a positive Perron vector. Normalise the vectors inside the compact probability simplex, then pass to a convergent subsequence.

\writeup Let $A\ge0$ be an arbitrary real $n\times n$ matrix. Write $J$ for the all-ones matrix and, for $\varepsilon>0$, put
\[
A_\varepsilon=A+\varepsilon J.
\]
Every entry of $A_\varepsilon$ is strictly positive, so Perron's theorem provides a real eigenvalue $\rho_\varepsilon>0$ and a vector $z_\varepsilon>0$ satisfying
\[
A_\varepsilon z_\varepsilon=\rho_\varepsilon z_\varepsilon.
\]
Scale $z_\varepsilon$ so that its coordinate sum is $1$. It then belongs to the compact simplex
\[
\Delta=\{x\in\R^n:x_i\ge0,\ \sum_i x_i=1\}.
\]

If $0<\varepsilon'<\varepsilon$, the entrywise inequality $A_{\varepsilon'}\le A_\varepsilon$ and Exercise~\ref{linear:ex:monotone} give
\[
\rho_{\varepsilon'}\le\rho_\varepsilon.
\]
The same comparison with $A$ gives $\rho(A)\le\rho_\varepsilon$. Thus as $\varepsilon\downarrow0$ the numbers $\rho_\varepsilon$ decrease to a finite limit $r\ge\rho(A)$.

Choose any sequence $\varepsilon_k\downarrow0$. Compactness supplies a subsequence, still denoted $\varepsilon_k$, such that $z_{\varepsilon_k}\to x\in\Delta$. The normalisation ensures $x\ne0$, although some of its coordinates may now be zero. Passing to the limit in the eigenvalue equation gives
\[
Ax=rx.
\]
Thus $r$ is an eigenvalue of $A$. By the definition of spectral radius, $r\le\rho(A)$. Together with the opposite inequality obtained above, this proves $r=\rho(A)$. Therefore
\[
\boxed{Ax=\rho(A)x,\quad x\ge0,\quad x\ne0}.
\]
Unlike the irreducible case, the limiting eigenvector need not have all coordinates positive.
\end{solution}

\begin{solution}{linear:ex:cycleperiod}
\idea Period is the gcd of lengths of all closed directed walks. A directed cycle has only multiples of its length. A connected undirected bipartite graph has a two-step return at each edge and no odd closed walk. Its bipartition also provides a sign-changing diagonal similarity.

\writeup For a directed cycle with $n$ vertices, the adjacency matrix $C$ cyclically shifts the basis vectors. A walk returns to its starting point precisely after a multiple of $n$ steps. The gcd of all positive return lengths is therefore $n$, so the period is $h=n$.

Let $\omega=e^{2\pi i/n}$. The vectors
\[
v_k=(1,\omega^k,\omega^{2k},\ldots,\omega^{(n-1)k})^{\mathsf T},
\qquad k=0,\ldots,n-1,
\]
are eigenvectors of the cyclic shift, with eigenvalues $\omega^k$ (or $\omega^{-k}$ if the shift is represented in the reverse convention). They form a Fourier basis, so the spectrum consists of the $n$ distinct $n$-th roots of unity. All have modulus $1=\rho(C)$, exactly matching the predicted peripheral spectrum.

Now let $A$ be the adjacency matrix of a connected undirected bipartite graph having at least one edge. Every step crosses from one part of the bipartition to the other, so every closed walk has even length. An edge traversed forward and then backward is a closed walk of length $2$. Therefore its period is $h=2$.

Let the two vertex classes be $U,V$, and define a diagonal matrix $D$ with entries $+1$ on $U$ and $-1$ on $V$. Every non-zero entry $a_{ij}$ joins opposite parts, so
\[
(DAD)_{ij}=D_{ii}a_{ij}D_{jj}=-a_{ij}.
\]
Since $D^{-1}=D$, this gives $DAD^{-1}=-A$. Thus $A$ and $-A$ are similar, and their spectra agree as multisets:
\[
\boxed{\lambda\in\sigma(A)\implies-\lambda\in\sigma(A)}
\]
with the same algebraic multiplicity. In particular the eigenvalues occur symmetrically about $0$.

To describe all of them, order the vertices by the two parts, of sizes $p,q$,
so $A=\bigl(\begin{smallmatrix}0&B\\B^{\mathsf T}&0\end{smallmatrix}\bigr)$.
The real symmetric matrix $BB^{\mathsf T}$ has an orthonormal eigenbasis by
the spectral theorem proved in Volume~1, Chapter~7. Its eigenvalues are
non-negative because $u^{\mathsf T}BB^{\mathsf T}u=\|B^{\mathsf T}u\|^2$.
Its kernel is exactly $\ker B^{\mathsf T}$, by the same identity, so it has
$r=\operatorname{rank}B$ positive eigenvalues, counting multiplicity.
For each eigenvector $u$ with eigenvalue $\mu>0$, direct multiplication gives
\[
A\binom{u}{\phantom{-}B^{\mathsf T}u/\sqrt\mu}
=\sqrt\mu\binom{u}{\phantom{-}B^{\mathsf T}u/\sqrt\mu},\qquad
A\binom{u}{-B^{\mathsf T}u/\sqrt\mu}
=-\sqrt\mu\binom{u}{-B^{\mathsf T}u/\sqrt\mu}.
\]
An orthonormal basis in each positive eigenspace yields independent vectors
in these two eigenspaces. The remaining eigenvectors are $(u,0)$ with
$u\in\ker B^{\mathsf T}$ and $(0,v)$ with $v\in\ker B$, all of eigenvalue $0$.
Their dimension is $(p-r)+(q-r)$. We have accounted for all $p+q$ dimensions:
the spectrum is the pairs $\pm\sqrt\mu$ and $p+q-2r$ zeros.
Different bipartite graphs give different $B$ and different numerical
eigenvalues. The excluded one-vertex graph has spectrum $\{0\}$ and no
positive closed-walk length, so no period under our definition.
\end{solution}

\begin{solution}{linear:ex:diagpos}
\idea A positive diagonal entry is a loop in the associated directed graph. A loop creates a closed walk of length $1$, forcing period $1$. If every vertex has a loop, a path can be padded to any longer desired length.

\writeup Let $A\ge0$ be irreducible of order $n$, and let $\Gamma(A)$ be its strongly connected directed graph. Suppose $a_{ii}>0$ for some $i$. Then there is a directed edge from $i$ to itself, a closed walk of length $1$. The period of an irreducible non-negative matrix is the gcd of the lengths of all closed walks. Since the set includes $1$, the gcd is $1$. By Theorem~\ref{linear:thm:primitive}, $A$ is primitive.

Now assume that \emph{all} diagonal entries are positive. Fix any ordered pair of vertices $(i,j)$. Strong connectivity provides a directed path from $i$ to $j$. Remove loops from that path until it is simple; a simple path in a graph with $n$ vertices uses at most $n-1$ edges. Let its length be $k$, so $0\le k\le n-1$.

At vertex $j$ we have a loop because $a_{jj}>0$. Follow the chosen path and then traverse this loop exactly $n-1-k$ times. The result is a directed walk from $i$ to $j$ of length exactly $n-1$. Every edge used has positive matrix weight, so this walk contributes a strictly positive term to the entry $(A^{n-1})_{ij}$.

As $i,j$ were arbitrary, every entry is positive:
\[
\boxed{A^{n-1}>0}.
\]
For $n=1$, interpret the zeroth power as $I_1>0$. The argument shows why diagonal loops greatly shorten the number of steps needed to make every entry positive.
\end{solution}

\begin{solution}{linear:ex:twostate}
\idea A two-state Markov chain has only one non-constant direction. The constant direction has eigenvalue $1$; the remaining eigenvalue is $1-a-b$. Splitting the identity into the two corresponding projections gives $P^k$ exactly.

\writeup Consider
\[
P=\begin{pmatrix}1-a&a\\b&1-b\end{pmatrix},
\qquad0<a,b<1.
\]
The stationary row vector $\pi=(\pi_1,\pi_2)$ satisfies $\pi P=\pi$ and $\pi_1+\pi_2=1$. The first coordinate equation is
\[
a\pi_1=b\pi_2.
\]
Solving with the normalisation gives
\[
\boxed{\pi=\left(\frac b{a+b},\frac a{a+b}\right)}.
\]

Since $P\mathbf1=\mathbf1$, one eigenvalue is $1$. The trace is $2-a-b$, so the other eigenvalue is
\[
\lambda=(2-a-b)-1=1-a-b.
\]
Our assumptions give $-1<\lambda<1$.

Define
\[
E=\mathbf1\pi=
\frac1{a+b}\begin{pmatrix}b&a\\b&a\end{pmatrix},
\quad
F=I-E=\frac1{a+b}\begin{pmatrix}a&-a\\-b&b\end{pmatrix}.
\]
These are complementary projections:
\[
E^2=E,\quad F^2=F,\quad EF=FE=0,\quad E+F=I.
\]
A direct multiplication gives $PE=E$ and $PF=\lambda F$, hence $P=E+\lambda F$. Raising this identity to the $k$-th power and using the projection relations yields
\[
\boxed{P^k=
\frac1{a+b}\begin{pmatrix}b&a\\b&a\end{pmatrix}
+\frac{(1-a-b)^k}{a+b}
\begin{pmatrix}a&-a\\-b&b\end{pmatrix}}
\qquad(k\ge1).
\]
At $k=0$ the power is $I$, separately; this avoids $0^0$ when $a+b=1$.
In that case $P=E$ and stationarity is reached after one step.
Because $|\lambda|<1$, the second term tends to zero and each row of $P^k$ tends to $\pi$. The convergence is geometric at rate
\[
\boxed{|1-a-b|^k}.
\]
It is fast when $a+b$ is near $1$ and slow when $a+b$ is near $0$ or $2$. The latter case alternates signs, corresponding to a chain close to deterministic switching.
\end{solution}

\begin{solution}{linear:ex:pagerankrate}
\idea The teleportation term sends every probability row to the same vector $v^{\mathsf T}$. It vanishes on the subspace of rows whose coordinates sum to zero. Every non-stationary left eigenvector belongs to that subspace.

\writeup Let $P$ be row stochastic, so $P\mathbf1=\mathbf1$. Let $v$ be a probability column vector, and suppose $0\le\alpha<1$. The Google matrix is
\[
G=\alpha P+(1-\alpha)\mathbf1v^{\mathsf T}.
\]
It is also row stochastic because $v^{\mathsf T}\mathbf1=1$, so $G\mathbf1=\mathbf1$.

Let $\lambda\ne1$ be an eigenvalue of $G$ and $x^{\mathsf T}\ne0$ a corresponding (possibly complex) left eigenvector:
\[
x^{\mathsf T}G=\lambda x^{\mathsf T}.
\]
Multiply on the right by $\mathbf1$. The stochastic identity gives
\[
x^{\mathsf T}\mathbf1
=x^{\mathsf T}G\mathbf1
=\lambda x^{\mathsf T}\mathbf1.
\]
Since $\lambda\ne1$, it follows that $x^{\mathsf T}\mathbf1=0$.

Therefore the rank-one teleportation term vanishes:
\[
x^{\mathsf T}G
=\alpha x^{\mathsf T}P
+(1-\alpha)(x^{\mathsf T}\mathbf1)v^{\mathsf T}
=\alpha x^{\mathsf T}P.
\]
For a complex row $y$, the triangle inequality and the row sums of $P$ give
$\|yP\|_1\le\|y\|_1$. Combining this with the eigenvalue equation yields
\[
|\lambda|\,\|x^{\mathsf T}\|_1
=\alpha\|x^{\mathsf T}P\|_1
\le\alpha\|x^{\mathsf T}\|_1.
\]
The eigenvector is non-zero, so division by its positive norm proves
$|\lambda|\le\alpha$, including $\alpha=0$.

For existence of a stationary row, start with any probability row $\mu_0$
and put $\mu_k=\mu_0G^k$. Differences of probability rows have sum zero,
so the same calculation gives
\[
\|\mu_{k+1}-\mu_k\|_1\le2\alpha^k,\qquad
\|\mu_{k+j}-\mu_k\|_1\le\frac{2\alpha^k}{1-\alpha}\quad(k\ge1).
\]
Each coordinate is Cauchy, hence converges by Volume~1, Chapter~3.
The limiting row $\pi$ is non-negative with sum $1$. Passing to the limit
in the finite multiplication $\mu_{k+1}=\mu_kG$ gives $\pi G=\pi$.
If $\pi'$ is another stationary probability row, then
$\|\pi-\pi'\|_1\le\alpha^k\|\pi-\pi'\|_1$ for every $k\ge1$,
so $\pi=\pi'$. Finally, iterating the zero-sum bound for $\mu-\pi$ gives
\[
\boxed{\|\mu G^k-\pi\|_1\le\alpha^k\|\mu-\pi\|_1\le2\alpha^k}
\qquad(k\ge1).
\]
For $\alpha=0$ every probability row is sent to $v^{\mathsf T}$ in one step.
No irreducibility assumption was needed, and $\pi$ may have zero coordinates.
\end{solution}

\begin{solution}{linear:ex:leslie}
\idea The survival rows determine an eigenvector from its first coordinate. Substituting that expression into the fertility row gives the Euler--Lotka equation. Its left side decreases strictly with a positive growth factor, so the positive root is unique.

\writeup Write the Leslie matrix as in Example~\ref{linear:exa:leslie}, with $f_i\ge0$, at least one $f_i>0$, and $s_i>0$. Define the survival products
\[
\ell_1=1,\qquad \ell_i=s_1s_2\cdots s_{i-1}\quad(i\ge2).
\]
All $\ell_i$ are positive. If $Lz=\lambda z$ with $\lambda\ne0$, rows $2$ through $n$ yield
\[
s_i z_i=\lambda z_{i+1},
\qquad
z_i=\ell_i\lambda^{-(i-1)}z_1.
\]
The first row then becomes
\[
\lambda z_1
=\sum_{i=1}^n f_i z_i
=z_1\sum_{i=1}^n f_i\ell_i\lambda^{-(i-1)}.
\]
Since a nonzero-eigenvalue eigenvector has $z_1\ne0$, divide by $\lambda z_1$ to get
\[
\boxed{\sum_{i=1}^n f_i\ell_i\lambda^{-i}=1}.
\]
After multiplication by $\lambda^n$, this is
\[
\lambda^n-\sum_{i=1}^n f_i\ell_i\lambda^{n-i}=0.
\]
To check multiplicities and zero roots as well, expand $\det(XI-L)$ along
the first row. The first term is $(X-f_1)X^{n-1}$. For $i\ge2$, the minor
of $-f_i$ forces the entries $-s_1,\ldots,-s_{i-1}$ and then $n-i$
diagonal entries $X$. The cofactor sign leaves $-f_i\ell_iX^{n-i}$.
Thus
\[
\boxed{\chi_L(X)=X^n-\sum_{i=1}^n f_i\ell_i X^{n-i}}.
\]

For existence of a positive root, let $m$ be the largest index with $f_m>0$.
The leading $m\times m$ Leslie block is irreducible: its survival edges
and the edge from its first row to its last column form a directed cycle
through all its vertices. For $m=1$ it is simply $[f_1]>0$.
Frobenius's theorem supplies a positive eigenvalue of this block.
The full matrix is block lower triangular, with a nilpotent trailing block,
so this is a positive root of $\chi_L$ as well.
For real $x>0$ put $g(x)=\sum_i f_i\ell_i x^{-i}$.
At least one coefficient is positive, so $g$ is strictly decreasing.
It follows that exactly one positive number $\rho$ satisfies $g(\rho)=1$.

The vector
\[
z=(1,\ell_2\rho^{-1},\ldots,\ell_n\rho^{-(n-1)})^{\mathsf T}>0
\]
satisfies $Lz=\rho z$. To identify the spectral radius without assuming
irreducibility of the full matrix, let $Ly=\lambda y$ and put
$t=\max_i|y_i|/z_i>0$. At a maximizing coordinate $i$,
$|\lambda|t z_i\le(L|y|)_i\le t(Lz)_i=t\rho z_i$.
Thus $|\lambda|\le\rho$ and the unique positive root is $\rho(L)$.
If all $f_i=0$ instead, $L$ is nilpotent, $\chi_L=X^n$ and there is no
positive root; this is precisely the case excluded from that assertion.

Finally, because $g$ is strictly decreasing and $g(\rho)=1$,
\[
\rho>1\iff g(1)>1
\iff\boxed{\sum_{i=1}^n f_i\ell_i>1}.
\]
The sum $\sum_i f_i\ell_i$ is the net reproductive rate. Equality to $1$ gives $\rho=1$, and a smaller rate gives $\rho<1$.
\end{solution}

\begin{solution}{linear:ex:doublystoch}
\idea A doubly stochastic matrix preserves the all-ones vector on both the left and the right. Its stationary distribution is therefore uniform. Primitivity makes all the other modes decay.

\writeup Let $P$ be an $n\times n$ primitive doubly stochastic matrix. Its row sums are $1$, so
\[
P\mathbf1=\mathbf1.
\]
Its column sums are also $1$, so
\[
\mathbf1^{\mathsf T}P=\mathbf1^{\mathsf T}.
\]
Choose the right Perron eigenvector $z=\mathbf1$ and the normalised left one $w^{\mathsf T}=\frac1n\mathbf1^{\mathsf T}$. Then $w^{\mathsf T}z=1$. The Perron--Frobenius convergence theorem for primitive matrices gives
\[
P^k\longrightarrow zw^{\mathsf T}
=\frac1n\mathbf1\mathbf1^{\mathsf T}
=\boxed{\frac1nJ}.
\]
Thus every row of $P^k$, which is a probability distribution after $k$ steps from one starting state, converges to the uniform vector.

Now suppose a fixed random card shuffle is applied repeatedly to a deck of $m$ cards, where there are $N=m!$ possible deck orders. Represent each order by a permutation $\sigma$; let the random shuffle permutation have law $\mu$. The transition probability from $\sigma$ to $\tau$ is, with a fixed composition convention,
\[
P_{\sigma,\tau}=\mu(\tau\sigma^{-1}).
\]
For fixed $\sigma$, the map $\tau\mapsto\tau\sigma^{-1}$ is bijective, so the row sum is $1$. For fixed $\tau$, the map $\sigma\mapsto\tau\sigma^{-1}$ is also bijective, so the column sum is $1$. Thus $P$ is doubly stochastic.

If the shuffle transition matrix is also primitive, the first part applies with $n=N=m!$:
\[
\boxed{P^k_{\sigma,\tau}\longrightarrow\frac1{m!}}
\]
for every starting order $\sigma$ and ending order $\tau$.
The primitivity hypothesis matters: on more than one deck order, a
deterministic shuffle has a permutation matrix and cannot converge to the
uniform law. The state-independent choice matters too: a positive transition
matrix with both rows $(3/4,1/4)$ is primitive but has a non-uniform stationary
law. It cannot arise from a fixed mixture of permutations of two states.
\end{solution}

\begin{solution}{linear:ex:wielandtsharp}
\idea The directed graph has one branching point, vertex $n$. From vertex $1$, every return must include the edge $n\to1$. Between returns we may make any number of shorter loops using the edge $n\to2$. This turns the problem into an integer equation.

\writeup Let $n\ge2$ and let $W$ be the adjacency matrix of the directed graph
\[
1\to2\to\cdots\to n,\qquad n\to1,\qquad n\to2.
\]
Every vertex $i<n$ has a single outgoing edge, leading to $i+1$. Therefore any non-empty closed walk based at vertex $1$ must first travel along
\[
1\to2\to\cdots\to n.
\]
At vertex $n$ there are only two choices. The edge $n\to1$ completes a circuit of length $n$ and returns us to the starting vertex. The edge $n\to2$ makes a loop
\[
n\to2\to3\to\cdots\to n
\]
of length $n-1$; after any number of these loops we can choose $n\to1$ to finish the circuit.

A closed walk at $1$ therefore consists of $a\ge1$ complete returns using the edge $n\to1$, together with $b\ge0$ shorter loops through $n\to2$, in some order. Its total length has the form
\[
\boxed{an+b(n-1),\qquad a\ge1,\ b\ge0}.
\]
Conversely, every such pair $(a,b)$ can be realised by taking $a$ full cycles and inserting $b$ short loops at a visit to $n$.

Suppose a closed walk had length $(n-1)^2$. We would have
\[
(n-1)^2=an+b(n-1).
\]
Reduce modulo $n-1$. Since $n\equiv1\pmod{n-1}$, this gives
\[
a\equiv0\pmod{n-1}.
\]
As $a\ge1$, we must have $a\ge n-1$. But then
\[
an\ge n(n-1)>(n-1)^2,
\]
contradicting the displayed equation. Therefore no such closed walk exists, and
\[
\boxed{(W^{(n-1)^2})_{11}=0}.
\]
The graph is strongly connected and has cycles of lengths $n$ and $n-1$, which are coprime. Hence $W$ is primitive. Theorem~\ref{linear:thm:wielandt} proves that $W^{(n-1)^2+1}>0$. The zero entry one step earlier shows that its exponent is \emph{exactly} $(n-1)^2+1$, proving the bound sharp. No numerical computation is needed.
\end{solution}

\section*{Sources and further reading}

\begin{itemize}
\item O.~Perron, \emph{Zur Theorie der Matrices}, Mathematische Annalen 64 (1907), 248--263; G.~Frobenius, \emph{\"Uber Matrizen aus
nicht negativen Elementen}, Sitzungsberichte der Preussischen Akademie (1912), 456--477.
\item L.~Collatz, \emph{Einschlie\ss ungssatz f\"ur die charakteristischen Zahlen von Matrizen}, Mathematische Zeitschrift 48
(1942), 221--226; H.~Wielandt, \emph{Unzerlegbare, nicht negative Matrizen}, Mathematische Zeitschrift 52 (1950), 642--648.
\item G.~Birkhoff, \emph{Extensions of Jentzsch's theorem}, Transactions of the AMS 85 (1957), 219--227.
\item H.~Schneider, \emph{Wielandt's proof of the exponent inequality for primitive nonnegative matrices}, Linear Algebra and its
Applications 353 (2002), 5--10 (with a facsimile, transcription and translation of the diary proof); A.~L.~Dulmage and
N.~S.~Mendelsohn, \emph{Gaps in the exponent set of primitive matrices}, Illinois Journal of Mathematics 8 (1964), 642--656.
\item A.~A.~Markov, \emph{Extension of the law of large numbers to dependent quantities} (in Russian), Izvestiya
Fiziko-Matematicheskogo Obshchestva Kazan 15 (1906), 135--156.
\item P.~H.~Leslie, \emph{On the use of matrices in certain population mathematics}, Biometrika 33 (1945), 183--212.
\item S.~Brin and L.~Page, \emph{The anatomy of a large-scale hypertextual web search engine}, Computer Networks and ISDN Systems
30 (1998), 107--117; K.~Bryan and T.~Leise, \emph{The \$25{,}000{,}000{,}000 eigenvector: the linear algebra behind Google},
SIAM Review 48 (2006), 569--581.
\item E.~Seneta, \emph{Non-negative Matrices and Markov Chains}, 2nd ed., Springer, 1981; C.~D.~Meyer, \emph{Matrix Analysis and
Applied Linear Algebra}, SIAM, 2000, Chapter~8, ``Perron--Frobenius Theory of Nonnegative Matrices.''
\end{itemize}
